%
% header_toggle_reconciliation.tex
%
% A Z specification of the lux collapsing_header open-state reconciliation:
% what the Display renders, and when it fires, across the window between a
% user's click on the disclosure triangle and the Hub's confirming re-push.
%
% Source of record:
%   src/punt_lux/display/renderers/imgui/collapsing_header.py
%       -- ImGuiCollapsingHeaderRenderer.begin: the honour + fire decision
%   src/punt_lux/display/renderers/imgui/header_open_arbiter.py
%       -- HeaderOpenArbiter: effective_open / should_fire, the pending slot
%   src/punt_lux/scene/widget_state.py
%       -- WidgetState: reset_honoured (the re-push reset),
%          discard_for (the removal clear)
%   src/punt_lux/scene/manager.py
%       -- SceneManager._replace_scene_state: drives reset_honoured on re-push
%
% Type-check:  fuzz -t docs/header_toggle_reconciliation.tex
% Model-check (see Section "Verification" for the exact goal predicates):
%   probcli docs/header_toggle_reconciliation.tex -model_check -p DEFAULT_SETSIZE 2
%   probcli docs/header_toggle_reconciliation.tex -model_check -nodead \
%       -goal "clicked=ztrue & steps>1" -p DEFAULT_SETSIZE 2
%   probcli docs/header_toggle_reconciliation.tex -model_check -nodead \
%       -goal "fires>1" -p DEFAULT_SETSIZE 2
%   probcli docs/header_toggle_reconciliation.tex -model_check -nodead \
%       -goal "spurious=ztrue" -p DEFAULT_SETSIZE 2
%   probcli docs/header_toggle_reconciliation.tex -model_check -nodead \
%       -goal "pending/=nopending & clicked=zfalse" -p DEFAULT_SETSIZE 2
%   probcli docs/header_toggle_reconciliation.tex -model_check -nodead \
%       -goal "rendered=ztrue & clicked=zfalse & shown/=hub" \
%       -p DEFAULT_SETSIZE 2
%
% Formatting follows the fuzz house rules: \quad~ for continuation
% indent (never \t1..\t9), schema lines under 80 columns, predicates
% broken at \land/\lor/\implies with the operator starting the
% continuation line.
%

\documentclass[a4paper,10pt,fleqn]{article}
\usepackage[margin=1in]{geometry}
\usepackage{fuzz}
\usepackage[colorlinks=true,linkcolor=blue,citecolor=blue,urlcolor=blue]{hyperref}

\begin{document}

\title{The lux Collapsing-Header Toggle Discipline: a Z Specification}
\author{Formal model of the render / fire / converge state machine}
\date{August 2026}
\maketitle

% ============================================================
\section{Introduction}
% ============================================================

A \texttt{collapsing\_header} carries one Hub-authoritative flag --- whether the
section is \emph{open} --- and is rendered every frame by an ImGui collapsing
header that keeps its \emph{own} stored open state.  Unlike a checkbox, whose
ImGui widget is handed the value as an argument and therefore renders exactly
what it is given, a collapsing header remembers.  \texttt{set\_next\_item\_open}
writes that memory; a click on the disclosure triangle flips it.  The renderer
must decide, every frame, what to write --- and that decision is the whole
subject of this document.

Writing the Hub value unconditionally on every frame is the obvious choice and
it is wrong.  It produces a \emph{double-step}: the click frame flips ImGui's
memory and the section opens; the very next frame writes the Hub value back
(the Hub has not yet heard about the click, so it still says closed) and the
section snaps shut; some frames later the Hub's confirming re-push arrives and
the section opens again.  Three rendered states, two of them transitions the
user did not ask for, all from one click that ends on the correct value.  This
is the defect the model exists to catch.

The correction is a single piece of per-render-session bookkeeping held in
\texttt{WidgetState}: the \emph{pending} value, the open state a
\texttt{HeaderToggled} has been fired for and the Hub has not yet ratified.
While a value is pending the renderer writes \emph{that} value rather than the
Hub's, so the click's own transition stands until the Hub speaks.  Every
Hub-driven operation clears the slot, which is what stops the optimism from
outliving its window and leaving the Display permanently disagreeing with the
Hub.

The same slot also suppresses re-firing.  The renderer fires when the state
ImGui reports differs from the state the renderer asked for, so once the click's
value is pending the following frames ask for it, are given it, and stay silent.
One slot serves both purposes, which is why this machine needs one where the
tab bar needs two: a tab bar compares against the Hub value to force-select and
against the Hub value again to fire, whereas a header compares against the
\emph{effective} value in both places.

This is a finite state machine over one header, so ProB can enumerate every
interleaving of render frames, clicks, re-pushes, and removals, and either
confirm the invariants or produce the exact offending trace.

% ============================================================
\section{Free Types and Constants}
% ============================================================

\begin{zed}
  ZBOOL ::= ztrue | zfalse
\end{zed}

A two-valued flag written as a free type rather than the toolkit
\texttt{bool}, so ProB animates it as a small enumerated set.  Every value in
this machine --- the Hub flag, ImGui's stored flag, the pending value --- is
one of these; a header's open state has exactly two possibilities, so unlike
the tab bar this model needs no basic type of identities at all.

\begin{zed}
  PENDING ::= nopending | held \ldata ZBOOL \rdata
\end{zed}

The optimistic slot either holds a value or does not.  Writing that as its own
free type rather than as a set of size at most one makes the empty and held
cases distinct by construction: there is no two-element value to constrain
away, and the state schema needs no structural predicate to rule one out.

\begin{zed}
  MAXFIRE == 2
\end{zed}

The fire counter saturates at \texttt{MAXFIRE}.  A single click may fire at
most once; a count of two is already a witness to the repeated-fire defect, so
counting no higher keeps the state space small without hiding the violation.

\begin{zed}
  MAXSTEP == 2
\end{zed}

The visible-transition counter saturates at \texttt{MAXSTEP}, for the same
reason.  One click may move the rendered state once; a count of two is the
double-step itself.

\begin{axdef}
  fresh : ZBOOL
\end{axdef}

\texttt{fresh} is the state ImGui gives a collapsing header it has never seen:
closed.  A newly declared header starts there whatever the Hub says, so the
first rendered frame must drive the declared value over this default --- the
cold-push case.

\begin{zed}
  fresh = zfalse
\end{zed}

% ============================================================
\section{State}
% ============================================================

The state is flat: one schema, eight components, and no predicate at all.  The
types carry what a predicate would otherwise have to assert, and the safety
properties (Section~\ref{sec:inv}) are deliberately \emph{absent} here, so that
a violating state remains reachable and ProB can find it.

\begin{schema}{Header}
  hub : ZBOOL \\
  shown : ZBOOL \\
  pending : PENDING \\
  clicked : ZBOOL \\
  rendered : ZBOOL \\
  fires : 0 \upto MAXFIRE \\
  steps : 0 \upto MAXSTEP \\
  spurious : ZBOOL
\end{schema}

\texttt{hub} is the Hub-authoritative open flag.  \texttt{shown} is ImGui's
stored open state --- what the user actually sees, since the widget draws its
arrow and gates its body from exactly this.  \texttt{pending} is the
optimistic slot: \texttt{nopending} when no toggle is outstanding, or
\texttt{held~v} for the value \texttt{v} a \texttt{HeaderToggled} has been
fired for and the Hub has not yet ratified.

The remaining four components are bookkeeping for the invariant checks, not
part of the renderer's own state.  \texttt{clicked} records whether a user
gesture is outstanding --- set by a click, cleared by any Hub-driven operation,
so it delimits exactly the click-to-re-push window.  \texttt{rendered} records
whether a frame has actually drawn the header since the last Hub-driven
operation; it is what lets invariant~5 say ``after a frame'' without a notion
of time.  \texttt{fires} counts fires within the current window, and
\texttt{steps} counts transitions of \texttt{shown} within it.
\texttt{spurious} latches \texttt{ztrue} the moment a fire occurs with no user
gesture behind it.

% ============================================================
\section{The Effective Value}
% ============================================================

Every frame renders the \emph{effective} value: the pending value while one is
outstanding, the Hub value otherwise.  It is a derived quantity, not state, so
it is given as a function of the two components it reads.

\begin{axdef}
  effective : ZBOOL \cross PENDING \fun ZBOOL
\where
  \forall h : ZBOOL @ effective(h, nopending) = h \\
  \forall h, v : ZBOOL @ effective(h, held~v) = v
\end{axdef}

The two cases exhaust \texttt{PENDING}, so the function is uniquely determined
--- there is one way to satisfy the axiom, and ProB has one constant to set up
rather than a family of them.

This one definition is the fix.  Replacing it with $effective(h, p) = h$ ---
the renderer writing the Hub value unconditionally --- is the buggy variant of
Section~\ref{sec:fidelity}, defect~(a), and it is what the shipped code did.

% ============================================================
\section{Initialization}
% ============================================================

\begin{schema}{Init}
  Header'
\where
  hub' = zfalse \\
  shown' = fresh \\
  pending' = nopending \\
  clicked' = zfalse \\
  rendered' = zfalse \\
  fires' = 0 \\
  steps' = 0 \\
  spurious' = zfalse
\end{schema}

A single initial state: a closed header, nothing pending, no frame drawn, no
fire seen.

% ============================================================
\section{Hub-Driven Operations}
% ============================================================

Three operations set the open flag from outside the renderer: the client
declaring the header, a whole-scene re-push of a surviving header, and a
removal followed by a re-added same-id header.  All three clear the pending
slot, because the Hub has now spoken and its value supersedes any optimism;
all three close the window (\texttt{clicked}, \texttt{fires}, \texttt{steps}
reset) and mark the header as not yet rendered under the new value.  That
clearing is the guard Section~\ref{sec:fidelity} removes to reproduce
defect~(b).

\texttt{Declare} models the client installing the header.  ImGui starts a fresh
widget, so \texttt{shown} is its closed default whatever the declared value is
--- the cold-push case invariant~5 covers.

\begin{schema}{Declare}
  \Delta Header \\
  v? : ZBOOL
\where
  hub' = v? \\
  shown' = fresh \\
  pending' = nopending \\
  clicked' = zfalse \\
  rendered' = zfalse \\
  fires' = 0 \\
  steps' = 0 \\
  spurious' = spurious
\end{schema}

\texttt{Repush} models a whole-scene re-push of a header that \emph{survives}
it (\texttt{SceneManager.\_replace\_scene\_state} calling
\texttt{reset\_honoured}).  The ImGui widget persists, so its stored
\texttt{shown} is kept; only the Hub value and the bookkeeping change.  This is
the operation the Hub performs to ratify a toggle --- and equally the one it
performs to \emph{reject} it, by re-pushing the value the user tried to move
away from.

\begin{schema}{Repush}
  \Delta Header \\
  v? : ZBOOL
\where
  hub' = v? \\
  shown' = shown \\
  pending' = nopending \\
  clicked' = zfalse \\
  rendered' = zfalse \\
  fires' = 0 \\
  steps' = 0 \\
  spurious' = spurious
\end{schema}

\texttt{RemoveReadd} models the header being removed and a same-id header
re-added (\texttt{discard\_for} clearing the slot).  Its state effect matches
\texttt{Repush} --- a stale ImGui memory kept under the same widget id, the
slot cleared --- but it travels a different code path, and each guards the
invariants on a path the other does not cover.

\begin{schema}{RemoveReadd}
  \Delta Header \\
  v? : ZBOOL
\where
  hub' = v? \\
  shown' = shown \\
  pending' = nopending \\
  clicked' = zfalse \\
  rendered' = zfalse \\
  fires' = 0 \\
  steps' = 0 \\
  spurious' = spurious
\end{schema}

% ============================================================
\section{The Render Frame}
% ============================================================

A frame either draws the header or does not.  A frame that draws writes the
effective value into ImGui's memory, lets the widget run, and reads back what
ImGui reports; the user may have clicked the disclosure triangle during that
run, in which case ImGui reports the opposite of what was written.  The two
drawing cases are therefore distinguished by whether the click happened, and
they share the write: in both, \texttt{shown} moves to the effective value
first.  A frame that does not draw is \texttt{FrameHidden}.

\texttt{FrameQuiet} is the no-click case.  ImGui is given the effective value
and reports it back unchanged, so nothing fires.  This single case covers the
first frame after a declare (driving the declared value over ImGui's closed
default), the steady state, the ratifying echo, and the convergence step after
a rejection --- in each, the frame simply shows what the renderer asked for.

\begin{schema}{FrameQuiet}
  \Delta Header
\where
  shown' = effective(hub, pending) \\
  steps' = \IF shown' = shown \THEN steps \\
\quad~  \ELSE (\IF steps < MAXSTEP \THEN steps + 1 \ELSE MAXSTEP) \\
  rendered' = ztrue \\
  hub' = hub \\
  pending' = pending \\
  clicked' = clicked \\
  fires' = fires \\
  spurious' = spurious
\end{schema}

\texttt{FrameClick} is the case where the user toggles the header during the
frame.  ImGui is given the effective value and the click flips it, so the
reported state is its opposite; that difference is what the renderer reads as a
genuine user toggle, and it fires exactly there.  The fired value goes into
\texttt{pending}, so the following frames ask ImGui for it and are given it ---
no further fire, and no revert.  A fire raised with no gesture behind it would
latch \texttt{spurious}; here the gesture is the operation itself, so
\texttt{clicked} is set.

\begin{schema}{FrameClick}
  \Delta Header
\where
  shown' \neq effective(hub, pending) \\
  pending' = held~shown' \\
  clicked' = ztrue \\
  rendered' = ztrue \\
  steps' = \IF shown' = shown \THEN 0 \ELSE 1 \\
  fires' = 1 \\
  hub' = hub \\
  spurious' = spurious
\end{schema}

The counters restart at the click: this frame opens a fresh window, in which
one fire and one visible transition are exactly what is allowed.  Writing
\texttt{steps' = 1} unconditionally would be wrong in one corner --- a click
that lands back on the state already shown moves no pixels --- so the
conditional keeps \texttt{steps} an honest count of transitions.

\texttt{FrameHidden} is the frame that does not draw: the header is nested in a
collapsed parent header or a hidden window, so its renderer never runs.
Nothing is written to ImGui's memory and nothing is read back, so every
component is untouched --- in particular \texttt{rendered}, which must not
claim a frame drew when none did.  It is enabled in \emph{every} state,
unguarded, modelling the fact that any frame may find the container collapsed.

\begin{schema}{FrameHidden}
  \Xi Header
\end{schema}

Because \texttt{FrameHidden} changes nothing it adds no reachable state; its
role is to make the not-drawn frame an explicit, always-available step.  With
\texttt{FrameQuiet} and \texttt{FrameClick} also unguarded, a render step is
enabled in every state, which is why the machine cannot deadlock.

% ============================================================
\section{Invariants}
\label{sec:inv}
% ============================================================

Five safety properties must hold across all reachable states, plus
deadlock-freedom.  They are stated here as predicates over \texttt{Header};
how each is discharged is given in Section~\ref{sec:verify}.  None is placed in
the \texttt{Header} schema: a predicate placed there would be enforced as a
guard on every successor, silently disabling any operation that would break it
and thereby masking the very defect this model exists to catch.

\paragraph{Invariant 1 --- one click, one visible step.}
Within the window a click opens, the rendered state moves at most once.  This
is the reported defect stated formally:
\[
  clicked = ztrue \implies steps \leq 1.
\]
The window is delimited by \texttt{clicked}, which the click sets and every
Hub-driven operation clears, so the invariant constrains exactly the frames
between the user's gesture and the Hub's answer.  Movement \emph{after} the Hub
answers is the Hub's prerogative --- a rejection legitimately moves the display
back --- and is governed by invariant~5 instead.

\paragraph{Invariant 2 --- no repeated fire.}
A single click yields at most one \texttt{HeaderToggled} across all frames
until the re-push.  With the counter reset on every click and re-push:
\[
  fires \leq 1.
\]

\paragraph{Invariant 3 --- no spurious fire.}
A \texttt{HeaderToggled} is fired only on a genuine user toggle, never off a
Hub echo, or a \texttt{fire $\to$ Hub $\to$ re-push $\to$ fire} loop would run:
\[
  spurious = zfalse.
\]

\paragraph{Invariant 4 --- no orphan optimism.}
The pending value is held only while a gesture is outstanding.  Every
Hub-driven operation clears it, so it can never outlive its window:
\[
  pending \neq nopending \implies clicked = ztrue.
\]
This is the invariant that bounds the optimism.  Without it the Display could
hold a value the Hub has contradicted for as long as the scene lives, which is
a worse failure than the flicker the pending slot was introduced to remove.

\paragraph{Invariant 5 --- the Hub wins once it has spoken.}
After any frame that drew, with no gesture outstanding, the rendered state is
the Hub's:
\[
  rendered = ztrue \land clicked = zfalse \implies shown = hub.
\]
This is Hub authority stated as a property of pixels.  It subsumes three cases
the design must get right: a header declared \emph{open} is shown open after
one frame even though ImGui's memory defaulted closed (the cold push); a
toggle the Hub \emph{ratifies} leaves the display on the ratified value; and a
toggle the Hub \emph{rejects} --- re-pushing the pre-click value --- pulls the
display back to it rather than stranding the user's optimistic view.

\paragraph{Invariant 6 --- deadlock freedom.}
No reachable state leaves the machine unable to progress.  All three frame
operations are unguarded, so a render step is enabled in every state.

% ============================================================
\section{Verification}
\label{sec:verify}
% ============================================================

Type-checking is a \texttt{fuzz} pass:
\begin{verbatim}
  fuzz -t docs/header_toggle_reconciliation.tex
\end{verbatim}

Invariants 1--5 are state properties, checked by asking ProB whether their
negation is \emph{reachable}.  A goal that is not found means the invariant
holds on every reachable state:
\begin{verbatim}
  # Invariant 1 (must report: goal NOT found)
  probcli docs/header_toggle_reconciliation.tex -model_check -nodead \
    -goal "clicked=ztrue & steps>1" -p DEFAULT_SETSIZE 2

  # Invariant 2 (must report: goal NOT found)
  probcli docs/header_toggle_reconciliation.tex -model_check -nodead \
    -goal "fires>1" -p DEFAULT_SETSIZE 2

  # Invariant 3 (must report: goal NOT found)
  probcli docs/header_toggle_reconciliation.tex -model_check -nodead \
    -goal "spurious=ztrue" -p DEFAULT_SETSIZE 2

  # Invariant 4 (must report: goal NOT found)
  probcli docs/header_toggle_reconciliation.tex -model_check -nodead \
    -goal "pending/=nopending & clicked=zfalse" -p DEFAULT_SETSIZE 2

  # Invariant 5 (must report: goal NOT found)
  probcli docs/header_toggle_reconciliation.tex -model_check -nodead \
    -goal "rendered=ztrue & clicked=zfalse & shown/=hub" \
    -p DEFAULT_SETSIZE 2
\end{verbatim}

Invariant 6 is the deadlock check over the full reachable state space:
\begin{verbatim}
  # Invariant 6 (must report: no deadlock, no invariant violation)
  probcli docs/header_toggle_reconciliation.tex -model_check \
    -p DEFAULT_SETSIZE 2
\end{verbatim}

% ============================================================
\section{Fidelity: reproducing the three defects}
\label{sec:fidelity}
% ============================================================

A model that cannot reproduce the known defect proves nothing.  Each defect
below is this specification with one guard removed, and each then makes the
corresponding invariant's negation reachable.  Defect~(a) is the shipped bug
this work fixes; (b) and (c) are the two ways a careless fix would go wrong,
included so the model constrains the design rather than merely describing it.
The three variants are the companion files
\texttt{header\_toggle\_reconciliation\_unconditional\_buggy.tex},
\texttt{header\_toggle\_reconciliation\_stale\_pending\_buggy.tex}, and
\texttt{header\_toggle\_reconciliation\_refire\_buggy.tex}.

\paragraph{Defect (a) --- the Hub value written unconditionally (the reported
double-step).}
Edit: $effective(h, p) = h$ for all $p$ --- the renderer calling
\texttt{set\_next\_item\_open(elem.open)} every frame, which is what the code
did.  Goal found: invariant~1 (\texttt{clicked=ztrue \& steps>1}), in ProB's
own shortest trace:
\begin{enumerate}
  \item \texttt{Init}: \texttt{hub = shown = zfalse}.
  \item \texttt{FrameClick}: the effective value is \texttt{hub = zfalse}, the
    click flips it, so \texttt{shown = ztrue} --- the section opens.
    \texttt{steps = 1}, \texttt{clicked = ztrue}, \texttt{pending = held~ztrue}
    (recorded, but the buggy \texttt{effective} ignores it).
  \item \texttt{FrameQuiet}: the next frame writes \texttt{effective = hub =
    zfalse}, so \texttt{shown = zfalse} --- the section snaps shut while the
    user's toggle is still in flight.  \texttt{steps = 2} with
    \texttt{clicked = ztrue}, violating invariant~1.
\end{enumerate}
A third step follows when \texttt{Repush(ztrue)} finally lands and the next
frame opens the section again, but the invariant has already been broken at
step~3: two transitions inside one window is the double-step the user sees.
Under the intact definition, step~3 writes the pending \texttt{ztrue}, so
\texttt{shown} does not move and \texttt{steps} stays 1.

\paragraph{Defect (b) --- the pending value surviving the re-push.}
Edit: \texttt{Repush} omits \texttt{pending' = nopending}, keeping
\texttt{pending' = pending} --- the optimistic hold not cleared when the Hub
speaks.  Goals found: invariant~4 (\texttt{pending/=nopending \&
clicked=zfalse}) and invariant~5 (\texttt{rendered=ztrue \& clicked=zfalse \&
shown/=hub}).  The invariant~5 trace is the one that shows the harm:
\begin{enumerate}
  \item \texttt{Init}; \texttt{FrameClick}: \texttt{shown = ztrue},
    \texttt{pending = held~ztrue}, \texttt{clicked = ztrue}.
  \item \texttt{Repush(zfalse)} in the buggy form --- the Hub \emph{rejects}
    the toggle, or an unrelated push carries the pre-click value:
    \texttt{hub = zfalse}, \texttt{clicked = zfalse}, but \texttt{pending}
    survives as \texttt{held~ztrue}.  This state alone violates invariant~4.
  \item \texttt{FrameQuiet}: the effective value is the stale
    pending \texttt{ztrue}, so \texttt{shown = ztrue} while
    \texttt{hub = zfalse} and no gesture is outstanding --- invariant~5
    violated.  Every later frame repeats it: the Display shows a value the Hub
    has contradicted, permanently.
\end{enumerate}
Under the intact reset, \texttt{Repush} empties \texttt{pending}, so the next
frame writes \texttt{hub} and the Display converges --- one visible step back,
which is a rejection correctly rendered.  ProB's own shortest witness for
invariant~4 reaches the same stale slot by a different route (two clicks, then
a re-push), confirming the defect is a property of the missing reset rather
than of one particular path to it.

\paragraph{Defect (c) --- firing against the Hub value rather than the
effective one.}
Edit: \texttt{FrameQuiet} fires when the state it rendered differs from
\texttt{hub}, instead of staying silent --- modelling a renderer that keeps the
pending slot for rendering but compares what ImGui reports against
\texttt{elem.open} when deciding to fire.  \texttt{FrameClick} is untouched, so
this is one guard, as the other two variants are.  Goal found: invariant~2
(\texttt{fires>1}), in ProB's shortest trace:
\begin{enumerate}
  \item \texttt{Declare(ztrue)}: a header declared open, ImGui still on its
    closed default.
  \item \texttt{FrameClick}: the frame writes the effective \texttt{ztrue} and
    the user's click flips it shut, firing once --- \texttt{fires = 1},
    \texttt{pending = held~zfalse}.
  \item \texttt{FrameQuiet}: the frame renders the pending \texttt{zfalse},
    which still differs from \texttt{hub = ztrue}, so it fires again ---
    \texttt{fires = 2}.  Every further frame in the window does the same: a
    fire storm for as long as the Hub takes to answer.
\end{enumerate}
Under the intact guard the comparison is against the effective value, which the
pending slot has already moved to the fired value, so the following frames find
no difference and stay silent.

\paragraph{Summary.}
Each defect's guard, removed, makes its invariant's negation reachable; with
every guard intact, all five reachability goals return \emph{not found} and the
deadlock check reports none.  That difference is the evidence that the model
detects the defect the code was changed to prevent, and that the pending slot
--- cleared by the Hub, compared against in both the render and the fire
decision --- is the minimal state that achieves it.

% ============================================================
\section{From the Model to the Code}
% ============================================================

The model's \texttt{pending} slot is one key in \texttt{WidgetState}, held by
\texttt{HeaderOpenArbiter}.  \texttt{effective} is the arbiter's
\texttt{effective\_open}, which the renderer passes to
\texttt{set\_next\_item\_open} in place of the raw \texttt{elem.open}; the
frame's fire decision is \texttt{should\_fire}, comparing the state
\texttt{collapsing\_header} reports against that same effective value and
recording the fired value into the slot.

The clearing of \texttt{pending} in \texttt{Repush} and \texttt{RemoveReadd} is
not new code: the slot takes one of \texttt{WidgetState}'s existing
per-render-session suffixes, which \texttt{reset\_honoured} wipes on every
scene replace and \texttt{discard\_for} wipes on removal.  That is the same
lifetime the tab bar's slots already have, and it is what invariant~4 depends
on --- the Display's optimism lasts exactly until the Hub next speaks.

\end{document}
